\documentclass[11pt]{article}
\usepackage[margin=25mm]{geometry}
\usepackage{fontspec}
\setmainfont{Latin Modern Roman}
\usepackage{amsmath,amssymb,mathtools}
\usepackage{unicode-math}
\setmathfont{Latin Modern Math}
\usepackage{xurl,hyperref}
\hypersetup{colorlinks=true,urlcolor=blue}
\setcounter{secnumdepth}{0}
\setlength{\parindent}{0pt}
\setlength{\parskip}{5pt}
\setlength{\emergencystretch}{3em}
\newcommand{\sref}[2]{\hyperlink{source:#1:#2}{[S#2]}}
\newcommand{\cataloguescope}[1]{\paragraph{Recorded target.}#1}
\newcommand{\R}{\mathbb R}
\newcommand{\C}{\mathbb C}
\newcommand{\N}{\mathbb N}
\newcommand{\Z}{\mathbb Z}
\newcommand{\Q}{\mathbb Q}
\newcommand{\F}{\mathbb F}
\newcommand{\eps}{\varepsilon}
\begin{document}
\section{Efficient Coarse-Grained Free-Energy Computation}\label{statement}

\subsection{Definitions and mathematical statement}

\paragraph{Shared notation.}$\mathbb N=\{1,2,\ldots\}$; the standard number fields have their usual meanings. All additional conventions are stated locally.


\paragraph{Definitions and conventions.} The atomistic coordinates split as $X=(X_r,X_c)$, with $N_r$ retained representative atoms and $N_c$ eliminated atoms. $\beta=(k_BT)^{-1}>0$ and $V(X_r,X_c)$ is the potential energy. The reduced free energy is the negative log of the conditional configurational integral with the retained coordinates fixed. A chain-like one-dimensional interaction geometry can still have individual atom positions in $\mathbb R^d$. The quantity $A(X_r)/N_c$ normalizes by the number of removed particles.

\paragraph{Source scope and qualifications.} The source already presents a computational strategy for one-dimensional nearest- and next-nearest-neighbor models using Markov-chain and transfer-operator arguments. The broad efficient-computation question lacks an error tolerance and complexity definition; it should not be advertised as having no known method in the source's chain case. This report passage is closely repeated at ID 30001271.

\paragraph{Catalogue mathematical target.}
\paragraph{Catalogue wording retained for traceability.}
For a one-dimensional atomistic system with repatom coordinates $X_r$ and eliminated coordinates $X_c$, define $A(X_r)=-\beta^{-1}\log\int\exp(-\beta V(X_r,X_c))\,dX_c$. Can $A(X_r)/N_c$ be computed efficiently as $N_c\to\infty$?

\paragraph{Saved literature context.} The supplied catalogue's assessment dated 2026-08-21 records: The question asks for an efficient large-Nc computation of the reduced-system free energy. The report supplies a formal expectation representation, but no general complexity or convergent computational theorem resolving the limit was found. The assessment is retained as dated provenance; it is not a new comprehensive literature or open-status certification.
\subsection{Short English statement}
Compute the free energy per eliminated particle efficiently in the thermodynamic limit of a chain-like atomistic system.
\subsection{Research attempt}
\paragraph{Model, date and investigated scope.}
OpenAI GPT-6 Astra (configured model \texttt{gpt-6-astra}), ultra reasoning; generated 11 October 2026. Computes harmonic-chain reduced free energy exactly and distinguishes fixed-length from fixed-strain limits.

\paragraph{Sources and literature check.}
The complete supplied definition, including source qualifications and the dated literature assessment, was reviewed. The calculation below is an elementary restricted attempt motivated by [S2]; it does not claim novelty or a fresh literature-status determination.

\paragraph{Context and investigated scope.} The complete supplied definition, including source qualifications and the dated literature assessment, was reviewed. The calculation below is an elementary restricted attempt motivated by [S2]; it does not claim novelty or a fresh literature-status determination.

\paragraph{An exactly soluble eliminated harmonic chain.} Fix endpoints $x_0=0,x_{n+1}=L$ and integrate $x_1,\ldots,x_n\in\mathbb R$ with $V=\frac k2\sum_{j=0}^n(x_{j+1}-x_j)^2$, $k>0$. Put $x_j=jL/(n+1)+y_j$ with $y_0=y_{n+1}=0$. The cross term telescopes, giving
\[
 V=\frac{kL^2}{2(n+1)}+\frac k2 y^TT_ny,
\]
where $T_n$ is tridiagonal with diagonal two and adjacent entries minus one. Its determinant obeys $d_n=2d_{n-1}-d_{n-2}$, $d_0=1,d_1=2$, hence $d_n=n+1$. Gaussian integration therefore yields
\[
 A_n=\frac{kL^2}{2(n+1)}-\frac n{2\beta}\log\frac{2\pi}{\beta k}+\frac1{2\beta}\log(n+1).
\]
This computes the conditional free energy without an $n$-dimensional numerical integral.

\paragraph{The scaling of retained coordinates matters.} For fixed $L$, $A_n/n\to-(2\beta)^{-1}\log(2\pi/(\beta k))$. If instead $L=(n+1)\ell$, the limit is $k\ell^2/2-(2\beta)^{-1}\log(2\pi/(\beta k))$. Thus the large-$N_c$ question needs a specification of endpoint strain. The reference length implicit in Lebesgue measure also fixes the additive free-energy normalization.

\paragraph{Remaining scope.} This is elementary Gaussian elimination for a nearest-neighbor quadratic potential, consistent with the source's already known chain strategies. It does not furnish uniform quadrature errors, mixing bounds or complexity for nonlinear nearest/next-nearest interactions. Extending the transfer-operator approximation with a stated accuracy criterion is the missing step. The closely repeated target at 30001271 is not treated as a distinct new theorem.

\paragraph{Verification and adversarial audit.}
Checked tridiagonal determinant recurrence, Gaussian normalization, n=1 and endpoint scaling; noted measure normalization and the repeated source passage. Model self-audit only; no independent human proof checking. Prompts A and B were applied to this canonical record and its qualified source scope. Checked tridiagonal determinant recurrence, Gaussian normalization, n=1 and endpoint scaling; noted measure normalization and the repeated source passage. Retained unresolved status and zero original-target progress; Certified approximation error and complexity for the source nonlinear chain interactions are missing.

Coordinator checked the telescoping quadratic shift, tridiagonal determinant recurrence, Gaussian normalization and both endpoint-scaling limits. This additional check was model-assisted, not independent human verification.

\paragraph{Tools, process and contributions.}
Python 3 was used for file preparation and structural checks; no numerical experiment or theorem prover was used. Contribution prompt A was followed by prompt B. GPT-6 Astra with ultra reasoning drafted the research and performed its stated self-audit; any additional coordinator contribution is identified in the verification and audit. A GPT-6-based Codex coordinator assembled the manuscript and checked repository integration; its exact variant and reasoning configuration were not exposed. No human mathematical contribution, independent human proof check or formal proof certification is supplied. Complete derivations and the consulted source references are included here; private reasoning transcripts are not published.

\paragraph{Final status and first remaining gap.}
\textbf{UNRESOLVED}. Certified approximation error and complexity for the source nonlinear chain interactions are missing.

\section{Completion Estimate}
\noindent\textbf{COMPLETION: 0\%}
The calculation is an auxiliary or restricted attempt and establishes no newly verified portion of the full source target. Certified approximation error and complexity for the source nonlinear chain interactions are missing.

\subsection{Sources}
\begin{itemize}
\item[\hypertarget{source:statement:1}{S1}] ULAM.AI and the UnsolvedMath Contributors, UnsolvedMath, catalogue entry. Catalogue curation under CC BY 4.0; definitions and English exposition added locally.
\url{https://huggingface.co/datasets/ulamai/UnsolvedMath}.
\item[\hypertarget{source:statement:2}{S2}] Interplay of Analysis and Probability in Physics (2008), Oberwolfach Reports, DOI 10.4171/owr/2008/55, pp. 3065--3138; official report source and mathematical context.
\url{https://doi.org/10.4171/owr/2008/55}.
\item[\hypertarget{source:statement:3}{S3}] Official EMS Press PDF of the same Oberwolfach report; PDF pages 9--10, printed pages 3073--3074.
\url{https://ems.press/content/serial-article-files/46198}.
\end{itemize}
\label{end:statement}
\end{document}
